\chapter{Returning to Discovery vs.\ Constitution}
\label{ch:resolving-fork}

\begin{plainlanguage}
\textbf{In plain terms:} Is intelligence discovered or constituted? This chapter closes that question — and it was actually closed quietly back in Chapter~5, not by any of the case studies in between. Some eliminations are forced by reality; some are pure rhetoric. Once you can tell them apart, intelligence can be neither purely discovered nor purely constituted — it's both, in different proportions for different parts of the concept.
\end{plainlanguage}

\section{The argument is already complete}

Chapter~\ref{ch:two-ways} opened with a fork left deliberately unresolved: is intelligence discovered, a fixed target that better theories progressively approximate, or constituted, nothing more than whatever the elimination process happens to converge on, with no further fact about correctness available? The fork was left open because, at that point in the book, no machinery existed that could answer it honestly. That machinery now exists, and the fork can be closed --- but it is worth being precise about exactly what closes it, because the temptation at this point in a book of this kind is to treat the accumulated weight of everything written since Chapter~\ref{ch:two-ways} as the argument, when in fact the argument was already complete by the end of Part~II, before any of the case studies were run.

The reasoning is this. Some eliminations are demonstrably world-forced: a mathematical counterexample, a physical hazard that reliably produces the harm it threatens, a system's measurable incapacity to perform a task. These score high on every component of the Chapter~\ref{ch:fidelity} fidelity functional simultaneously, independent of who asserts them. Other eliminations are demonstrably rhetorical: amplified, low-directness, low-closure, authority-dependent --- the empty bowl, the moral panic. Distinguishing the two requires the fidelity machinery; treating them identically, as either pure discovery (everything is world-forced, nothing is rhetorical) or pure constitution (everything is rhetorical, nothing is world-forced) would assert what the fidelity components were built specifically to deny. Therefore intelligence --- or any concept subject to this kind of perspectival elimination --- can be neither entirely discovered nor entirely constituted. Some of what survives elimination survives because the world will not permit otherwise. Some of what survives does so only because no one has yet supplied, or no culture has yet been willing to supply, the constraint that would eliminate it.

Parts~III through V of this book do not establish this argument. They stress-test it: Part~III asks whether the underlying elimination dynamics behave sensibly in messy, real cases (Chapter~\ref{ch:room-for-river}, Chapter~\ref{ch:development}); Part~IV asks whether a real cultural pedagogy already organizes itself around something like the fidelity distinction this book is formalizing (Chapter~\ref{ch:counterfoil} through Chapter~\ref{ch:recursion}); Part~V asks whether the recursive compression conjecture generalizes beyond intelligence to a second competence word (Chapter~\ref{ch:creativity}), and what would explain competence-word persistence if the conjecture is right (Chapter~\ref{ch:why-compress}). None of these tests, however they turn out, can overturn the Chapter~\ref{ch:two-ways}--\ref{ch:fidelity} argument, because that argument does not depend on any of them succeeding --- and Part~V is, in fact, the clearest available demonstration of this independence, since it did not turn out cleanly. Chapter~\ref{ch:creativity} found only partial transport: the recognition/implementation shape recurs in creativity research, but creativity needed additional axes (domain expertise, social uptake) the intelligence case never surfaced, complicating rather than confirming the conjecture's specific two-axis form. Chapter~\ref{ch:why-compress} left its central conjecture explicitly unfinished, with no working cost model. Neither outcome touches the weak-realist resolution stated in this chapter. What narrows, given Part~V's actual result, is only the generality claim of Chapter~\ref{ch:recursion}'s conjecture --- which is now better stated as ``recursive compression is not unique to intelligence'' than as ``recursive compression takes the same shape everywhere,'' a real but more modest finding than a clean second success would have produced.

\section{Stating the position with its actual content}

\begin{definition}[Weak-Realist Decomposition]
Restate Chapter~\ref{ch:two-ways}'s hard/soft split without the appearance of a formal limit it has not earned:
\[
\Theta = H \cup S, \qquad H = \left\{\theta \in \Theta \;:\; \begin{array}{l}\theta\text{ survives arbitrarily extended}\\ \text{sequences of high-fidelity elimination}\end{array}\right\}.
\]
\end{definition}

This is deliberately less formal-looking than an earlier version of this definition, which expressed $H$ as a limit, $\lim_{t\to\infty}$, of a set-valued process. That phrasing was set aside because a limit of a set-valued sequence requires specifying what converges and in what topology --- apparatus this book has not built and does not need. The set-builder form above says exactly what is meant and no more: $H$ is whatever, as a matter of fact, keeps surviving no matter how long and how varied a sequence of high-fidelity perspectival contact is run against it. $S$ is everything else --- the dimensions that have moved, or could plausibly move, under different cultural or historical contact. Stating it this way is a small change, but it matters for the same reason the rest of this book has insisted on small changes of this kind throughout: the notation should never look more proven than the argument behind it.

\begin{conjecture}[Persistence, demoted from ``Theorem'']
Elements of $H$ remain invariant under what might be called \emph{cultural reparameterization} --- a transformation of the local soft constraints $S$ that a culture or era applies, while $H$ is held fixed.
\end{conjecture}

This cannot honestly be stated as a theorem, because the reparameterization operation it claims invariance under has not been formally defined anywhere in this book. Doing so would require specifying precisely what counts as a transformation of a culture's constraint set, and on what grounds $H$ should be expected to be a fixed point of that transformation rather than merely a slow-moving one. Left as a conjecture, it restates the hard-core intuition in slightly more technical language without proving anything new about it --- which is the honest status to assign it given what has actually been derived.

\section{What this resolves and what it does not}

The hard core of any competence bundle --- the dimensions forced by contact with real, resistant failure --- looks discovered, stable across cultures and eras because it is not up to cultures and eras to revise it. The soft periphery --- caricature, rhetorical amplification, locally constituted emphasis --- looks constituted, and shifts, because nothing forces it to stop. ``Intelligence'' names the union of both, and this is exactly why arguments about its definition have never fully resolved across a century of serious attempts: part of the disagreement is real, unsettled empirical territory about where the boundary of $H$ actually sits, and part of it is a contest, often conducted by people who believe themselves to be doing the former, over which locally-constituted constraints in $S$ get to count, rhetorically, as though they belonged to $H$. The discovery cycle that opened Chapter~\ref{ch:discovery-trouble} is, on this account, not a series of failures to find the right definition. It is the visible record of that contest being fought out, repeatedly, across a century in which the tools to tell the two sides of the fight apart did not yet exist in stated form.

\section{A Positive Criterion: Convergent Corroboration}
\label{sec:corroboration}

\begin{plainlanguage}
\textbf{In plain terms:} So far this book only knows how to rule things out. But two completely independent fields sometimes arrive at the same structural claim with no contact between them — that's active confirmation, not mere survival, and deserves its own name. The worked case: a 1930s linguistic theory, confirmed decades later by neuroscience that had no idea the theory existed.
\end{plainlanguage}

\subsection{A case where theory arrived before the evidence that confirmed it}

In the mid-1930s, Roman Jakobson and Nikolai Trubetzkoy proposed that phonemes --- the basic sound units linguists had long treated as atomic --- are not atomic at all. Each phoneme, on their account, decomposes into a small bundle of abstract binary distinctions: voicing, place of articulation, and several others. This was a purely theoretical claim, argued on structuralist linguistic grounds, with no access to and no stake in any future neuroscience. There was, at the time, no method available that could have looked inside a living brain and checked.

Roughly seventy years later, intracranial electrode recordings taken from patients undergoing pre-surgical mapping for epilepsy --- a technique unrelated in origin or motivation to 1930s phonology --- found something specific: electrodes in auditory cortex do not specialize for individual phonemes. They specialize for exactly the kind of abstract distinctive features Jakobson and Trubetzkoy had proposed on paper, decades before anyone could test it. As Yosef Grodzinsky describes the finding in a 2026 lecture discussing the case, the coincidence is striking enough to call ``sci-fi'': a theoretical construct from one discipline, developed with no neuroscientific input, turned out to be the literal organizing scheme of a physical organ, confirmed by a second discipline that had no reason to vindicate the first.

This deserves to be stated narrowly and precisely, because it is tempting --- and would be a fidelity violation in exactly the sense Chapter~\ref{ch:fidelity} warns against --- to inflate the claim beyond what the evidence supports. What is established here is the auditory-cortex electrode finding itself: a specific structural prediction, made on theoretical grounds, confirmed by a specific later study using an independent methodology. It is not established, by this case alone, that distinctive feature theory has been continuously validated across phonology, acquisition, speech-error research, and aphasiology for a century --- that broader claim may well be true, but this book has not verified it, and citing it as already demonstrated here would be borrowing exactly the kind of unearned narrative force this book's own machinery exists to catch.

What this case offers the book is not a confirmed grand history. It is a single, clean, well-documented instance of something the existing formalism has no name for.

\subsection{Survival is not corroboration}

Everything built in Chapters~\ref{ch:candidate-space} through \ref{ch:fidelity} is negative machinery. A constraint $C_i$ rules candidates \emph{out}; a candidate's standing in $H$ is, so far, defined entirely by what it has not yet been eliminated by. Call this \emph{survival}: $\mu_t(\theta)$ stays bounded away from $0$ because nothing has yet driven it down.

The distinctive-features case is not a survival story. Nobody eliminated a rival candidate and left distinctive feature theory standing by default. Something stronger happened: a second, independent methodology, with no access to and no stake in the first theory, arrived at the same structural claim on its own. Call this \emph{corroboration}: $\mu_t(\theta)$ rises, actively, because independent perspectives keep recovering it rather than merely failing to rule it out.

These are different processes, and the difference matters for what kind of confidence each one licenses. A candidate that has merely survived could still be wrong in a way no one has yet found the right perspective to detect --- its standing is a statement about the limits of the search so far, not about the candidate itself. A candidate that keeps being independently recovered by methodologies that could not have influenced one another is in a different epistemic position: the explanation that it is simply wrong has to also explain why unrelated investigative routes, with no contact between them, kept arriving at it anyway. That is a much harder coincidence to construct.

This motivates the section's central definition, stated before any of the supporting machinery beneath it:

\begin{definition}[Convergent Corroboration]
A candidate $\theta$ is \emph{convergently corroborated} when it remains recoverable across observational regimes --- independently re-derived by perspectives whose mutual independence is high, especially when the regimes are separated by enough methodological or temporal distance that the later one could not have been shaped, even unconsciously, by the earlier one.
\end{definition}

The Jakobson--Trubetzkoy case is the concrete instance motivating this definition, and it is worth being precise about which fact is doing the work. It is not merely that structural phonology and intracranial electrophysiology both, as it happens, endorse the same distinction. It is that the two methodologies are about as different as two ways of investigating language could be --- and that the gap between them is large enough, roughly seventy years, that the second methodology did not yet exist when the first theory was proposed. Distinctive feature theory could not have been reverse-engineered from a foreknowledge of how electrode recordings would later turn out, because no one running the 1930s analysis had any such recordings to consult, and no one running the later study needed the 1930s theory to design their experiment around. The independence is not merely high; in this specific respect it is closer to structurally guaranteed.

\subsection{Why independence is the construct the section actually depends on}

The Lagrangian apparatus built later in this section is, on reflection, not where this section's real claim lives. The claim lives entirely in how independence is defined and measured, and it is worth stating plainly that the whole argument for convergent corroboration as a meaningful category stands or falls on this one construct.

\begin{definition}[Independence]
For two perspectives $P_i, P_j$, let
\[
\mathrm{Ind}(i,j) \in [0,1]
\]
denote one minus the fraction of shared evidentiary or methodological lineage between them, estimated from documented history: shared authors, shared instruments, shared funding sources, shared theoretical ancestry, or direct citation contact. $\mathrm{Ind}(i,j)$ is maximal when the two perspectives developed with no contact and no shared assumptions.
\end{definition}

Two further, more specific conditions push $\mathrm{Ind}(i,j)$ toward its maximum, and both hold in the case motivating this section. \emph{Methodological distinctness}: the two perspectives investigate the candidate through evidence of different kinds entirely --- structural analysis of phonological patterning versus electrical activity recorded directly from cortical tissue, sharing no instruments, no data, and no inferential machinery. \emph{Temporal impossibility of contamination}: when the confirming methodology did not yet exist at the time the original claim was made, the confirming study cannot have been designed, even informally, around producing the expected result --- there is no mechanism by which the earlier theory could have biased the later measurement, because the later measurement's entire technical apparatus postdates the theory it ends up confirming. This second condition is doing more work than it might first appear: it is close to a structural analogue of a pre-registered prediction, except stronger, since pre-registration only rules out a researcher consciously shaping a study around a known hypothesis, while temporal impossibility rules out the possibility of any such shaping occurring at all, conscious or not.

\subsection*{Strong versus weak independence}

The discussion above gives two named conditions that push $\mathrm{Ind}(i,j)$ toward its maximum, but a reader trying to apply this machinery to a new case needs more than two extreme examples to work from. Is developmental psychology independent of linguistics? Is a second neuroscience laboratory, using the same scanner technology as the first, independent of it? These are exactly the questions a working taxonomy needs to answer, and a three-tier version, built from the two factors already in play --- shared theoretical ancestry and shared instrumentation --- is enough to organize most cases that will actually arise.

\emph{Weak independence}: shared theory and shared instruments. Two laboratories using the same scanner technology, trained in the same theoretical tradition, and citing each other's prior work sit here. Agreement between them is closer to consensus than corroboration in the sense of Section~\ref{sec:corroboration}'s later subsection on that distinction --- informative, but only weakly so, because a shared assumption or a shared instrumental artifact could produce the agreement independently of whether the underlying candidate is real.

\emph{Moderate independence}: shared theory, different instruments. Two research programs that accept the same broad theoretical framework but reach their conclusions through different measurement techniques --- behavioral testing versus neuroimaging, for instance, within a shared cognitive-science paradigm --- rule out instrumental artifact as a common cause of agreement, but not theoretical bias: both could still be finding what their shared framework predisposed them to look for.

\emph{Strong independence}: different theory, different instruments. This is the tier the Jakobson--Trubetzkoy case occupies, and it is the tier that does the most epistemic work, because it rules out both common causes of false agreement at once --- neither a shared instrument nor a shared theoretical expectation can explain why two frameworks with nothing in common arrived at the same structure.

This taxonomy is offered as a practical aid to applying $\mathrm{Ind}(i,j)$ rather than as a further formal refinement of it --- the three tiers are not claimed to exhaust the space or to be sharply bounded from one another, and real cases will often sit ambiguously between them. What the taxonomy is meant to prevent is the more common error of treating any agreement between two named, separately-credentialed fields as automatically strong independence, when the relevant question is always the finer-grained one: how much theory and instrumentation do the two perspectives in question actually share, not merely what their fields are conventionally called.

\subsection*{Temporal independence}

One further factor is worth separating out from the theory/instrument taxonomy above, because it behaves differently from either axis. The temporal impossibility of contamination condition already given is binary: either the confirming methodology existed at the time of the original claim, making some form of shaping conceivable, or it did not, ruling shaping out entirely. But time separation does additional work beyond this binary gate, and that additional work is graded rather than binary.

A recovery event fifty or seventy years after the original proposal is stronger corroborating evidence than a recovery event one week later, even holding theoretical and instrumental independence fixed, because longer separation makes alternative explanations for the agreement progressively less plausible: less chance of an unrecorded personal connection between the researchers involved, less chance that the original claim was still active enough in a field's working memory to bias a later study's design or interpretation, more opportunity for the original claim to have been quietly forgotten or fallen out of use, only to be rediscovered rather than confirmed. The Jakobson--Trubetzkoy case is, again, closer to the strong end of this than a typical example: roughly seventy years separate the original proposal from the confirming electrode recordings, a gap long enough that the original theoretical claim was not part of any live research conversation the electrophysiologists were responding to. Call this \emph{temporal independence}, a graded factor that strengthens $\mathrm{Ind}(i,j)$ continuously as the gap widens, distinct from and additional to the binary impossibility-of-contamination condition: a recovery can satisfy the binary condition (the methodology genuinely did not exist yet) while still being temporally weak in this graded sense (the gap was short enough that the claim was still circulating actively when the confirmation arrived), and the strongest cases --- this book's worked example among them --- satisfy both.

Without a properly demanding definition of independence, convergent corroboration collapses into nothing more than consensus --- two perspectives agreeing because they share assumptions, instruments, or community, which is a much weaker and more familiar phenomenon, fully susceptible to the same borrowed-authority failure mode Chapter~\ref{ch:bad-counterfoils} describes (an entire field can be wrong together, in agreement, for reasons having nothing to do with the world). With a properly demanding definition, the same surface pattern --- multiple perspectives endorsing the same candidate --- becomes something categorically different: rediscovery.

\subsection{Recovery, alongside elimination}

The existing machinery already has a quantity tracking how strongly perspectives argue \emph{against} a candidate at a given perspective: $e_t(\theta)$, aggregated across perspectives as elimination pressure in the instability potential below. Convergent corroboration needs the positive counterpart.

\begin{definition}[Endorsement]
For a perspective $P_i$ and a candidate $\theta \in \Theta$, let
\[
\varepsilon_i(\theta) \in [0,1]
\]
denote the degree to which $P_i$'s own framework, evaluated on its own terms and stated prior to any corroborating evidence, specifically predicts $\theta$ --- as opposed to merely permitting $\theta$ among many candidates the framework happens not to forbid.
\end{definition}

Distinctive feature theory scores high on this measure with respect to the auditory-cortex finding: Jakobson and Trubetzkoy's framework did not merely fail to forbid feature-based neural organization. It specifically predicted that any physical encoding of phonemic contrast, wherever it was eventually found, would be organized by abstract feature rather than by raw acoustic identity --- a falsifiable structural commitment made decades before the technology to check it existed.

\begin{definition}[Recovery Count]
\[
R(\theta,t) = \bigl|\{(i,j) : i<j\le n(t),\ \varepsilon_i(\theta) > \tau,\ \varepsilon_j(\theta) > \tau,\ \mathrm{Ind}(i,j) > \tau'\}\bigr|
\]
for fixed thresholds $\tau,\tau' \in (0,1)$ --- the number of sufficiently independent perspective pairs that each, on their own terms, specifically endorse $\theta$.
\end{definition}

$R(\theta,t)$ is the direct positive counterpart to an elimination count: where elimination tallies how many perspectives a candidate has survived contact with, recovery tallies how many sufficiently independent perspectives have actively required it. A candidate with $R(\theta,t) = 0$ may still be entirely defensible --- most candidates this book treats as part of $H$ have never been tested for corroboration at all, only for survival, and the distinction is not meant to demote them. What $R(\theta,t) > 0$ adds is a second, stronger kind of standing that survival alone cannot confer.

With endorsement and independence both in hand, the weighted version follows:

\begin{definition}[Pairwise and Aggregate Corroboration]
\[
K_{ij}(\theta) = \varepsilon_i(\theta)\,\varepsilon_j(\theta)\,\mathrm{Ind}(i,j), \qquad K(\theta,t) = \sum_{i<j \le n(t)} K_{ij}(\theta).
\]
\end{definition}

This is genuinely computable, not merely evocative, and it is worth being explicit about what computing it would actually involve rather than leaving the claim abstract: take the phoneme clustering implied by distinctive feature theory and the phoneme clustering implied by the electrode-response data, and measure their agreement above chance --- an adjusted Rand index or normalized mutual information between the two independently-derived partitions would serve. That number is what $K(\theta)$ is measuring in this specific case. It has not been computed here; the underlying electrode data was not available to this book. The procedure is real and stated precisely enough to be carried out by someone with access to the data. The number itself remains an open empirical question, and the chapter should not be read as implying otherwise.

\subsection{A proposed dynamical model}

The definitions above --- survival versus corroboration, independence, recovery count, and weighted corroboration --- are where this section's actual content lives, and the Jakobson--Trubetzkoy case motivates and supports every one of them directly. What follows is different in kind and explicitly subordinate to it: a heuristic dynamical model, offered as a formal description of a process whose existence has already been established by the case above, not as the source of the section's authority. A reader who finds the Lagrangian unconvincing or unnecessary loses nothing of the section's substantive claim, which is fully stated already.

\begin{optionalformal}
The operational definitions above are the part of this section the book can fully stand behind. What follows is different in kind: a heuristic dynamical model, built so that it reduces, in a stated limit, to machinery the book has already established. It earns inclusion only to the extent that reduction holds, and the reduction is offered here as a conjecture, not a completed derivation.

Treat $\mu_t(\theta)$ as a continuum field over $\Theta$, using the conceptual-distance metric $d_\Theta$ from Chapter~\ref{ch:discovery-trouble} to give nearby candidates meaning. Define the \emph{instability potential}
\[
V(\theta,t) = \underbrace{\sum_i w_i(t)\,e_i(\theta)}_{\text{elimination pressure}} \;-\; \lambda\, K(\theta,t), \qquad \lambda > 0.
\]
Elimination pressure raises $V$, pushing plausibility down. Corroboration lowers $V$, carving a stable well a candidate can settle into.

\begin{definition}[Belief-Field Lagrangian, proposed dynamical model]
\[
\mathcal{L}[\mu,\dot\mu,\nabla_\theta\mu] = \frac{\kappa}{2}\dot\mu(\theta,t)^2 \;-\; \frac{\sigma}{2}\bigl(\nabla_\theta \mu(\theta,t)\bigr)^2 \;-\; V(\theta,t)\,\mu(\theta,t).
\]
\end{definition}

The middle term is the one genuinely new structural addition here, and it formalizes something the book has so far only used informally: the discovery cycle of Chapter~\ref{ch:discovery-trouble} narrows toward \emph{neighborhoods} of definitions, not isolated points, which presupposes exactly the kind of smoothness this term enforces --- two candidates close in conceptual distance should not differ wildly in plausibility absent specific evidence distinguishing them.

The associated Euler--Lagrange equation is
\[
\kappa\,\ddot\mu(\theta,t) - \sigma\,\nabla^2_\theta \mu(\theta,t) + V(\theta,t) = 0.
\]

\begin{conjecture}[Overdamped Reduction]
In the regime where belief revision is slow relative to inertial effects ($\kappa\ddot\mu \approx 0$), the dynamics above reduce to gradient flow,
\[
\sigma\,\nabla^2_\theta \mu(\theta,t) \approx V(\theta,t),
\]
and a discrete-time, renormalized version of this flow recovers something structurally close to the corrected soft-elimination update of Chapter~\ref{ch:fidelity}: $\mu$ decreases where elimination pressure dominates ($V > 0$), increases or stabilizes where corroboration dominates ($V < 0$), and diffuses smoothly across nearby candidates rather than updating each $\theta$ in isolation.
\end{conjecture}

This is stated as a conjecture because it is plausible rather than demonstrated: the actual discretization and renormalization needed to confirm that this continuum model reproduces Chapter~\ref{ch:fidelity}'s update exactly, rather than merely resembling it qualitatively, has not been carried out. The gap is left visible deliberately, in keeping with the rest of this book's practice, rather than smoothed over by the elegance of the notation surrounding it.

\begin{definition}[Conjugate Momentum and Hamiltonian]
With $\pi(\theta,t) = \partial\mathcal{L}/\partial\dot\mu = \kappa\dot\mu(\theta,t)$, the Legendre transform gives
\[
\mathcal{H}[\mu,\pi] = \int_\Theta d\theta \left[\frac{\pi(\theta,t)^2}{2\kappa} + \frac{\sigma}{2}\bigl(\nabla_\theta\mu(\theta,t)\bigr)^2 + V(\theta,t)\,\mu(\theta,t)\right].
\]
\end{definition}

This is mechanically correct given the Lagrangian above --- a genuine Legendre transform, not a separately invented expression --- but its physical interpretation should not be overstated. Nothing in this book provides independent evidence that belief revision has well-defined momentum or conserves $\mathcal{H}$ in any literal sense. The Hamiltonian is included because it follows directly from the proposed Lagrangian and because the gradient-flow limit connecting back to Chapter~\ref{ch:fidelity} is the part doing real work; the energy-conservation language that comes bundled with Hamiltonian mechanics is borrowed for formal convenience, not asserted as a discovered property of how concepts actually move.
\end{optionalformal}

\subsection{Corroboration is not agreement}

One further confusion is worth heading off explicitly, since it is the most natural way to misread everything built in this section: convergent corroboration is not consensus, and the two should not be allowed to blur together.

Consensus concerns how many people, or how many restatements, endorse a claim. Corroboration concerns how many \emph{independent routes} recover it. Three researchers who each repeat the same textbook argument may increase social confidence in a claim, but contribute little new epistemic information, because all three endorsements trace back to a single common source --- in the terms of this section, $\mathrm{Ind}(i,j) \approx 0$ for any pair among them, and the resulting agreement, however widely repeated, adds almost nothing to $K(\theta,t)$. Three genuinely independent methodologies that each, on their own terms and without contact with one another, recover the same candidate are in a different epistemic position entirely, because $\mathrm{Ind}(i,j) \approx 1$ for each pair, and the probability of all three converging on the same structure through entirely unrelated routes, if that structure were not real, is correspondingly lower.

This is precisely why convergent corroboration measures something popularity cannot. A claim repeated by thousands of people can have low corroboration in this book's sense if every repetition traces back to the same original source --- the moral panic of Chapter~\ref{ch:bad-counterfoils} is, among other things, an instance of exactly this: wide agreement, near-zero independence. A claim supported by only a handful of perspectives can have high corroboration if those perspectives reached it by genuinely unrelated paths.

The case motivating this section illustrates the distinction precisely, and it is worth restating its scope narrowly rather than letting the point above invite overreach. What this book has established is convergence between two regimes specifically: 1930s structural phonology and 2010s intracranial electrophysiology, with $\mathrm{Ind}(i,j)$ close to maximal between that pair for the documented reasons already given. It would be a natural next step to suggest that a broader convergence --- across language acquisition research, speech pathology, and other adjacent fields, in addition to the two regimes already discussed --- would make the corroboration case stronger still, and that may well be true. But this book has not verified that broader convergence, and asserting it here, however tempting given how well it would fit the argument, would repeat exactly the kind of unearned extension this section's own earlier subsection warned against. The two-perspective case, stated at the scope it has actually been demonstrated, is already a strong and sufficient illustration of what corroboration measures that consensus alone cannot: not the size of a coalition, but the independence of the paths by which a structure is recovered.

\subsection{What this section establishes and what it leaves open}

The substantive claim of this section is now fully stated before any physics notation appears: survival and corroboration are different processes, convergent corroboration is well-defined once independence is given a sufficiently demanding operational meaning, and the Jakobson--Trubetzkoy case is a real, correctly attributed instance of it --- with $\mathrm{Ind}(i,j)$ close to maximal for documentable reasons (methodological distinctness, temporal impossibility of contamination), and $R(\theta,t) > 0$ for that case on any reasonable choice of threshold. None of that depends on the dynamical layer that follows. The Lagrangian, its overdamped reduction, and the Hamiltonian are offered explicitly as a secondary formal description of a phenomenon the section has already established by other means --- a conjecture-level model, built to connect to existing machinery in a specific limit, with that connection not yet rigorously confirmed. Distinguishing these two layers matters more here than almost anywhere else in the book, because the temptation to let the elegance of a Lagrangian formulation lend unearned credibility to the operational claims beneath it is exactly the kind of borrowed authority Chapter~\ref{ch:bad-counterfoils} was built to diagnose. A reader who finds the dynamical model unpersuasive loses nothing of the section's actual argument; a reader who finds it persuasive should be clear that its persuasiveness is borrowed from the case and the definitions that precede it, not the other way around.
